\documentclass[10pt,a4paper]{article}
\usepackage[utf8]{inputenc}
\usepackage[T1]{fontenc}
\usepackage[margin=1.8cm]{geometry}
\usepackage{enumitem}
\usepackage{booktabs}
\usepackage{parskip}
\setlength{\parskip}{0.35em}

\title{\textbf{AI-Powered Smart Agriculture Advisory System for Smallholder Farmers in Uganda}\\
\large Project Planning and Execution Readiness (February Deliverables)}
\author{Project Team}
\date{}

\begin{document}
\maketitle

\section*{Three Core Deliverables}
The project is organized into three core deliverables that build toward the final output:
\begin{enumerate}[noitemsep]
  \item \textbf{11th February:} Improved functional prototype and trained disease model baseline.
  \item \textbf{28th February:} Integrated advisory system with analytics and preliminary results.
  \item \textbf{Final stage:} Consolidated system with validation evidence for wider deployment.
\end{enumerate}

This report focuses on the first two deliverables as requested.

\section*{Deliverable 1 (11th February)}
\textbf{Brief description:} An improved prototype where core user flows (disease detection, weather, and basic advisory) are functional end-to-end.

\textbf{1) Objective and research alignment}
\begin{itemize}[noitemsep]
  \item Directly addresses disease diagnosis from uploaded images.
  \item Addresses weather-informed recommendations for farmer decisions.
\end{itemize}

\textbf{2) Key components/features}
\begin{itemize}[noitemsep]
  \item Disease module: image upload, preprocessing, ONNX inference, treatment mapping.
  \item Weather module: real weather integration and advisory generation.
  \item Frontend pages and authentication workflow for usable access.
\end{itemize}

\textbf{3) Technical work involved}
\begin{itemize}[noitemsep]
  \item Prepare image dataset and class mapping (\texttt{label\_map.json}).
  \item Split data (train/validation/test), train model, export ONNX.
  \item Integrate backend APIs with frontend and run end-to-end checks.
\end{itemize}

\textbf{4) Scope and limitations}
\begin{itemize}[noitemsep]
  \item Scope: stable functionality for supported model classes/crops.
  \item Limitation: model depends on available training classes and image quality.
\end{itemize}

\textbf{5) TRL:} \textbf{TRL 5} (validated in a relevant development environment).

\section*{Deliverable 2 (28th February)}
\textbf{Brief description:} A more complete integrated system with analytical modules and preliminary performance evidence.

\textbf{1) Objective and research alignment}
\begin{itemize}[noitemsep]
  \item Extends decision support by combining disease, weather, market, and advisory knowledge.
  \item Produces preliminary evidence of technical feasibility and reliability.
\end{itemize}

\textbf{2) Key components/features}
\begin{itemize}[noitemsep]
  \item Integrated advisory: disease + weather + seasonal + RAG guidance.
  \item Market module: realistic regional reference prices and trend/prediction endpoints.
  \item Reliability controls: confidence/margin/entropy logic and crop-gated checks.
\end{itemize}

\textbf{3) Technical work involved}
\begin{itemize}[noitemsep]
  \item Refine inference handling for real-world images and uncertain cases.
  \item Standardize weather responses (including wind speed and units).
  \item Integrate market endpoints and frontend visualizations.
  \item Document preprocessing and implementation flow for reproducibility.
\end{itemize}

\textbf{4) Scope and limitations}
\begin{itemize}[noitemsep]
  \item Scope: integrated and testable advisory platform with preliminary outputs.
  \item Limitation: not yet a full field pilot; continued calibration is needed with more real user images.
\end{itemize}

\textbf{5) TRL:} \textbf{TRL 6} (demonstrated as an integrated system in a relevant environment).

\section*{How These Deliverables Build to the Final Outcome}
\begin{itemize}[noitemsep]
  \item Deliverable 1 proves core technical feasibility.
  \item Deliverable 2 proves integrated decision-support capability.
  \item The final deliverable builds on both to produce a stronger validated system.
\end{itemize}

\begin{table}[h]
\centering
\small
\begin{tabular}{@{}lll@{}}
\toprule
\textbf{Item} & \textbf{11th February} & \textbf{28th February} \\
\midrule
Type & Functional prototype + trained model baseline & Integrated system + preliminary analytical results \\
Main objective & Disease + weather feasibility & Multi-module advisory integration and reliability \\
TRL & 5 & 6 \\
\bottomrule
\end{tabular}
\end{table}

\end{document}
